\chapter{GIỚI THIỆU TỔNG QUAN}
\label{ch:gioi_thieu}

\section{Đặt vấn đề}

Thị giác máy tính ba chiều đang là nền tảng cho hàng loạt ứng dụng thực tế: cánh tay robot gắp
vật trong kho, hệ thống kiểm kê tự động, thiết bị hỗ trợ người khiếm thị, hay các ứng dụng thực
tại tăng cường cần biết đồ vật đặt ở đâu và lớn cỡ nào. Điểm chung của những ứng dụng này là
chúng không cần một nhãn phân lớp đẹp, mà cần một mô tả \emph{định lượng} về cảnh: có bao nhiêu
vật, mỗi vật chiếm khoảng không gian nào, và các vật quan hệ với nhau ra sao.

Khoảng cách giữa yêu cầu nêu trên và năng lực hiện có của hệ thống thị giác nằm ở chỗ dữ liệu đầu vào chỉ là một mảng giá trị điểm ảnh, trong đó khái niệm "một vật thể" không được biểu diễn tường minh. Việc khôi phục khái niệm này, cùng với các thuộc tính định lượng đi kèm, bộc lộ một số điểm nghẽn kỹ thuật cụ thể.

\begin{enumerate}[leftmargin=1.5cm] 
\item \textbf{Ảnh màu hai chiều đã đánh mất thông tin chiều sâu.} Từ một tấm ảnh RGB đơn thuần
không thể suy ra kích thước thật theo centimet, cũng không biết vật nào đứng trước vật nào. Cảm
biến RGB-D bổ sung được thông tin này, nhưng chỉ ghi nhận được bề mặt nhìn thấy, và trả về giá
trị sai trên các vật trong suốt, phản chiếu hoặc màu tối \cite{suchi2019ocid}.

\item \textbf{Tập chủng loại đồ vật trong nhà gần như vô hạn.} Một bộ phân lớp đóng huấn luyện
trên $N$ lớp cố định sẽ thất bại ngay khi gặp vật ngoài danh sách, và tệ hơn, nó vẫn trả về một nhãn với độ tự tin cao, vì hàm softmax trên tập lớp đóng buộc phải phân bổ toàn bộ khối lượng xác suất cho các lớp có sẵn và không có cơ chế biểu diễn trạng thái từ chối. Cách tiếp cận class-agnostic, gộp mọi vật tiền cảnh về một lớp duy nhất object, loại bỏ được hạn chế này nhưng đánh mất hoàn toàn thông tin danh tính. \cite{xie2021uois}.

\item \textbf{Che khuất là quy luật chứ không phải ngoại lệ.} Trong cảnh lộn xộn, phần lớn đồ
vật bị che một phần. Cảm biến chỉ đo được phần nhìn thấy, nên mọi con số kích thước đều là ước
lượng của \emph{phần hiện hữu} chứ không phải của toàn bộ vật thể. Nhập nhằng giữa hai đại lượng
này là nguồn sai lệch hệ thống thường bị bỏ qua khi báo cáo kết quả.
\end{enumerate}

Nút thắt thứ hai gợi ra một ý tưởng trung tâm cho đề tài: nếu gọi được tên vật, ta lập tức truy
cập được một kho tri thức tiên nghiệm về kích thước thật của lớp vật đó. Tri thức này có thể
dùng để kiểm tra xem phép đo có hợp lý không, và xa hơn, để \emph{hiệu chỉnh} phép đo khi độ sâu
đầu vào không có đơn vị --- trường hợp luôn xảy ra với các mô hình ước lượng độ sâu đơn mắt
\cite{ranftl2020midas}. Đề tài đặt ra ba câu hỏi nghiên cứu tương ứng.

\noindent\textbf{Câu hỏi nghiên cứu 1.1} \textit{Việc bổ sung thông tin độ sâu vào bộ lọc có
thực sự cải thiện độ chính xác đếm đối tượng so với phương án chỉ dùng ảnh màu hay không?}

\noindent\textbf{Câu hỏi nghiên cứu 1.2} \textit{Một mô hình nhận diện từ vựng mở có cơ chế từ
chối đạt được độ chính xác bao nhiêu trên ảnh đời thường, và khoảng cách giữa độ chính xác trên
mặt nạ đã khớp với độ chính xác đầu-cuối lớn đến mức nào?}

\noindent\textbf{Câu hỏi nghiên cứu 1.3} \textit{Cơ chế hiệu chỉnh tỉ lệ toàn ảnh dựa trên tri
thức kích thước tiên nghiệm có làm giảm sai số đo trên dữ liệu giữ lại hay không, khi được đánh
giá bằng thủ tục để-lại-một để loại bỏ tính vòng tròn?}

\section{Mục tiêu đề tài}

Đề tài hướng đến các mục tiêu chính sau:

\begin{itemize}[leftmargin=1.5cm]
\item Xây dựng bộ phân vùng instance class-agnostic hoạt động ổn định trên cảnh bàn lộn xộn có
che khuất, huấn luyện và đánh giá trên dữ liệu RGB-D thật.
\item Thiết kế bộ đếm đối tượng hai tầng và định lượng phần đóng góp riêng của thông tin độ sâu
bằng một thí nghiệm loại trừ (ablation).
\item Bổ sung tầng nhận diện từ vựng mở có cơ chế từ chối, tách rời khỏi bộ phân vùng, với từ
vựng ràng buộc theo bảng tri thức kích thước.
\item Ước lượng kích thước ba chiều của phần nhìn thấy từ mặt nạ và độ sâu, kèm cơ chế kiểm tra
tính hợp lý so với dải kích thước thật của lớp vật.
\item Đề xuất và \emph{kiểm chứng nghiêm ngặt} cơ chế hiệu chỉnh tỉ lệ toàn ảnh, bao gồm cả thí
nghiệm kiểm soát độ nhạy lẫn thủ tục đánh giá để-lại-một.
\item Suy luận quan hệ không gian theo luật và đóng gói toàn bộ kết quả thành đồ thị cảnh dạng
JSON có thể tái sử dụng.
\end{itemize}

\section{Đối tượng và phạm vi nghiên cứu}

\textbf{Đối tượng nghiên cứu:} các cảnh bàn và cảnh sàn nhà chứa nhiều đồ gia dụng đặt lộn xộn,
có che khuất lẫn nhau, được ghi nhận bằng cảm biến RGB-D đã căn chỉnh; cùng với các ảnh màu đời
thường khi cần đánh giá riêng phần nhận diện.

\textbf{Phạm vi nghiên cứu.} Đề tài đo \emph{phần nhìn thấy} của mỗi vật, không đặt mục tiêu
khôi phục hình học bị che khuất; mọi con số kích thước vì vậy là ước lượng xấp xỉ chứ không phải
phép đo toàn vật. Việc phân vùng được giữ class-agnostic với đúng một lớp tiền cảnh; danh tính
được gán ở tầng suy luận tách rời và bị giới hạn trong 87 lớp mà bảng tri thức có dải
kích thước tương ứng. Phần quan hệ không gian chỉ xét các cặp hai vật với tập vị từ hữu hạn, dựa
trên luật hình học chứ không học từ dữ liệu. Đề tài không xử lý chuỗi thời gian, không theo vết
đối tượng qua các khung hình, và không tối ưu tốc độ suy luận cho thiết bị nhúng.

Một giới hạn cần nêu rõ ngay: các dải kích thước tiên nghiệm là khoảng giá trị điển hình của cả
lớp vật, không phải nhãn thật của từng cá thể. Do đó cờ báo tính hợp lý chỉ có giá trị bắt các
sai lệch thô, chứ không phải một chỉ số sai số chính xác.

\section{Đóng góp của đề tài}

\begin{enumerate}[leftmargin=1.5cm]
\item Một quy trình xử lý tám bước hoàn chỉnh, chạy được đầu-cuối, tích hợp phân vùng
class-agnostic, đếm có lọc theo độ sâu, nhận diện từ vựng mở, đo ba chiều, kiểm tra tính hợp lý,
hiệu chỉnh tỉ lệ và suy luận quan hệ; kèm 74 kiểm thử tự động và một ứng dụng web minh hoạ.

\item Một thiết kế ràng buộc giữa từ vựng nhận diện và bảng tri thức kích thước, bảo đảm bộ nhận diện chỉ phát ra những nhãn mà hệ thống có tri thức kích thước tương ứng, nhờ đó hai thành phần không thể sai lệch với nhau khi hệ thống được mở rộng.

\item Một thủ tục đánh giá hiệu chỉnh tỉ lệ phá vỡ tính vòng tròn, dùng dữ liệu công khai và
không cần gán nhãn thủ công: mỗi vật được đo hai lần từ cùng khung hình với cùng mặt nạ thật,
chỉ thay đổi nguồn độ sâu, nên chất lượng phân vùng bị loại khỏi vai trò biến gây nhiễu.

\item \textbf{Một kết quả âm tính được đo đúng cách và báo cáo đầy đủ}: cơ chế hiệu chỉnh tỉ lệ
vượt qua thí nghiệm kiểm soát độ nhạy nhưng vẫn làm sai số đo tăng 24,2\% trên tập dữ liệu giữ
lại. Đề tài phân tích nguyên nhân thay vì lược bỏ kết quả này, và chỉ ra rằng việc đạt bài kiểm
soát là điều kiện cần chứ không phải điều kiện đủ.
\end{enumerate}

\section{Bố cục báo cáo}

Phần còn lại của báo cáo được tổ chức như sau. Chương 2 trình bày cơ sở hình học của ảnh RGB-D, các kiến trúc phân vùng và nhận diện được sử dụng, tính bất định về tỉ lệ của độ sâu đơn mắt, và phân tích vấn đề đánh giá vòng tròn vốn là nền tảng cho thiết kế thực nghiệm của đề tài. Chương 3 mô tả ba nguồn dữ liệu, các bước tiền xử lý, kiến trúc hệ thống tám bước và bộ chỉ số đánh giá, trong đó trọng tâm là thủ tục loại-một-ra dùng cho phần hiệu chỉnh tỉ lệ. Chương 4 trình bày và bàn luận kết quả của sáu thí nghiệm. Chương 5 đối chiếu kết quả với mục tiêu, nêu hạn chế và đề xuất hướng phát triển.
